% LaTeX preamble for the assembled manual. Loaded by build-manual.R.

% pdflscape rotates BOTH the page content and the page in the PDF viewer, so a
% landscape table is read by turning the screen rather than the head. lscape
% alone rotates only the content. This is what page-fit.lua wraps a table in on
% the rare occasion one is too wide for a portrait page.
\usepackage{pdflscape}

% Tables that run past one page continue with their header repeated rather than
% being silently truncated, which is what happens to a plain tabular.
\usepackage{longtable}
\usepackage{booktabs}

% Wrapping inside a cell. Without it a long cell sets on one line and pushes the
% table past the page edge. page-fit.lua decides how wide each column should be;
% this is what lets a cell use the width it was given rather than run past it.
\usepackage{array}
\usepackage{ragged2e}

% The hex sticker's fill, for links and chapter rules. Sampled from
% man/figures/logo.png rather than guessed.
\usepackage{xcolor}
\definecolor{RavenOrange}{HTML}{EC7414}

% Code blocks wrap instead of running off the page. The docker commands in the
% usage article are long, and a reader cannot copy what was never printed.
\usepackage{fvextra}
\DefineVerbatimEnvironment{Highlighting}{Verbatim}{
  breaklines, breakanywhere, commandchars=\\\{\}, fontsize=\small}

% Chapter openings carry the brand rule rather than the default plain heading.
\usepackage{titlesec}
\titleformat{\chapter}[display]
  {\normalfont\huge\bfseries}
  {\color{RavenOrange}\chaptertitlename\ \thechapter}
  {12pt}{\Huge}
\titlespacing*{\chapter}{0pt}{0pt}{20pt}

% Running header naming the chapter, so a printed copy says which article a page
% came from.
\usepackage{fancyhdr}
\pagestyle{fancy}
\fancyhf{}
\fancyhead[L]{\footnotesize\nouppercase{\leftmark}}
\fancyhead[R]{\footnotesize cyRAVEN}
\fancyfoot[C]{\footnotesize\thepage}
\renewcommand{\headrulewidth}{0.4pt}

% A landscape page keeps the same footer, otherwise the page number rotates off
% the visible area.
\fancypagestyle{lscape}{\fancyhf{}\fancyfoot[C]{\footnotesize\thepage}}

% Images wider than the text block are scaled down rather than overrunning it.
% The figures in the worked example are full-width panel grids.
%
% THE HEIGHT CAP IS WHAT KEEPS A FIGURE WITH ITS TITLE. pandoc bounds an image
% to the full text height, and a figure exactly as tall as the page leaves
% nowhere for the heading and the line above it: they stay on the page before and
% the figure is on its own, which is the split this cap exists to prevent. At
% 0.82 there is room left for the heading, the line describing it, the caption
% and the space a float leaves around itself. Verified on the built manual: all
% 31 figures that have a heading of their own keep it on their own page.
%
% It binds only on a figure taller than it is wide, which here is six of
% forty-seven: the panel grids are wider than they are tall and are limited by
% the line width, not by this.
\usepackage{graphicx}
\setkeys{Gin}{width=\linewidth, height=0.82\textheight, keepaspectratio}

% Figures are printed where they are written rather than floated to wherever
% they next fit. These are documentation figures discussed by the sentence above
% them, not plates that can be read in any order, and a float that moves is a
% figure whose text is now on another page.
\usepackage{float}
\floatplacement{figure}{H}

% RESERVING THE PAGE A TABLE OR FIGURE AND ITS TITLE NEED.
%
% \cyneed{<height>} starts a new page when less than that height is left on this
% one, so a heading and the line under it move to the page their table or figure
% lands on rather than being stranded above the break. page-fit.lua emits one
% call before each heading, sized from the title's own length.
%
% WHY NOT THE needspace PACKAGE. \needspace asks for the space with stretchable
% glue and a penalty of -100, which is a breakpoint TeX may CONSIDER later, once
% more material has been contributed. When the material after it is a longtable,
% later means after the table has begun, and longtable reacts to a break inside
% its own extent by repeating its header on the next page. The manual came out
% with a stray column header floating above the heading of the table it belonged
% to. The comparison here is made against \pagegoal at the moment of the call and
% the break, if one is needed, is \newpage, whose penalty of -10000 fires the
% output routine on the spot, before any of the table exists.
%
% \pagegoal is \maxdimen until something has been contributed to the page. That
% means the page is empty, so there is nothing to break away from.
\newcommand{\cyneed}[1]{%
  \par
  \begingroup
  \dimen0  = \dimexpr#1\relax
  % Clamped to one page: a request larger than a page can never be satisfied, so
  % honouring it would eject a blank page and change nothing.
  \ifdim\dimen0 >\textheight \dimen0 =\textheight \fi
  \ifdim\pagegoal=\maxdimen\else
    \ifdim\dimexpr\pagegoal-\pagetotal\relax<\dimen0  \newpage \fi
  \fi
  \endgroup}

% A figure's height is known only once the file has been opened, so the image is
% measured into a box that is never printed and the reservation is made from
% that. #2 is the number of lines the heading and description above it take.
%
% \IfFileExists rather than assuming: the vignettes guard every include_graphics
% with file.exists(), so a missing illustration is meant to leave a gap in the
% text, not stop the build.
\newsavebox{\cyfigbox}
\newcommand{\cyfigneed}[2]{%
  \IfFileExists{#1}{%
    \sbox\cyfigbox{\includegraphics{#1}}%
    \cyneed{\ht\cyfigbox+\dp\cyfigbox+#2\baselineskip}%
  }{}}

% URLs break at any character. Several in the documentation are long enough that
% \sloppy alone leaves them protruding into the margin.
\usepackage{xurl}
